%%%%%
%
\section{Vitesse de rotation � partir
 de la d�riv�e des angles d'attitude}


\itemstyle{Op\'erateur~:} \operfmt{Vrot\_repav}
\index{Vrot\_repav}

%%%%%OBJET%%%%
\itemstyle{Objet~:}
Cet op�rateur calcule les vitesses de rotations autour des axes d'une
plate-forme mobile � partir de la d�riv�e des angles d'attitude.

%%%%%ARGUMENTS%%%%
\bigskip
%%%%%ENTREES%%%%
\begin{tabarg}
\multicolumn{3}{l}{\titretab {Entr\'ees}}\\
\hline
$\phi$ & Vecteur, \texttt{float} & Angle de roulis, ou g�te,
       positif lorsque le cot� tribord est enfonc� ($rad$)\\
$\theta$ & Vecteur, \texttt{float} & Angle de tangage, ou assiette
       longitudinale, positif lorsque l'avant est relev� ($rad$)\\
$d\phi$ & Vecteur, \texttt{float} & D�riv�e temporelle du roulis ($rad/s$)\\
$d\theta$ & Vecteur, \texttt{float} & D�riv�e temporelle du tangage ($rad/s$)\\
$d\psi$ & Vecteur, \texttt{float} & D�riv�e temporelle du cap ($rad/s$)\\
\end{tabarg}

\medskip
%%%%%SORTIES%%%%
\begin{tabarg}
\multicolumn{3}{l}{\titretab {Sorties}}\\
\hline
$Vr_x$ & Vecteur, \texttt{float} & Vitesse de rotation autour de 
l'axe longitudinal ($rad/s$)\\
$Vr_y$ & Vecteur, \texttt{float} & Vitesse de rotation autour de 
l'axe transversal ($rad/s$)\\
$Vr_z$ & Vecteur, \texttt{float} & Vitesse de rotation autour de
l'axe vertical  ($rad/s$)\\
\end{tabarg}


%%%%%FORMULE%%%%
\itemstyle{Formule}

%%Vrot_repav
\begin{eqnarray*}
Vr_x & = & d\phi - \sin \theta . d\psi \\
Vr_y & = & \cos \phi .d\theta + \sin \phi .\cos \theta .d\psi \\
Vr_z & = & - \sin \phi .d\theta + \cos \phi .\cos \theta .d\psi \\
\end{eqnarray*}


\itemstyle{Commentaire :}

Le rep�re li� � la plate-forme est d�fini de la mani�re suivante :
\begin{itemize}
\item axe X : longitudinal � la plate-forme, orient� vers l'avant
\item axe Y : transversal � la plate-forme, orient� vers tribord
\item axe Z : vertical � la plate-forme, orient� vers le fond
\end{itemize}

%%%%%AUTEUR%%%%
\itemstyle{Auteur~:}
CNRM/GMEI/TRAMM

%%%%%REFERENCE%%%%

\itemstyle{R\'ef\'erences~:}

Bulletin NCAR-RAF Nr 23 \cite{RAF-23},
norme AFNOR NF-X-02-115 "\textit{Symboles et vocabulaire
de la m�canique du vol}"\cite{AFNOR-74}, page 7.

